\section{Registerkartan}
\label{sec:regmap}

Tretton register, vart och ett 32 bitar. \textbf{Indexet} är det som går ut på ledningen;
offsetkolumnen är alltid $\text{index} \times 4$ och är ett arv från en äldre, minnesmappad variant
av samma konstruktion.

\begin{booktable}{@{}rlllL@{}}
\thead{Idx} & \thead{Register} & \thead{Offset} & \thead{Åtk.} & \thead{Beskrivning} \\ \midrule
0 & \code{STATUS} & \code{0x00} & R & Bit 0: TX klar. Bit 1: RX giltig. Bit 2: Fel. \\
1 & \code{TX_ID} & \code{0x04} & R/W & 11-bitars sändar-ID (bitarna 10--0). \\
2 & \code{TX_DLC} & \code{0x08} & R/W & Data Length Code (bitarna 3--0, värde 0--8). \\
3 & \code{TX_DATA_LO} & \code{0x0C} & R/W & Sändardata byte 0--3 (byte 0 = MSB). \\
4 & \code{TX_DATA_HI} & \code{0x10} & R/W & Sändardata byte 4--7. \\
5 & \code{TX_SEND} & \code{0x14} & W & Skriv \code{0x1} för att utlösa sändning. \\
6 & \code{RX_ID} & \code{0x18} & R & Mottagen identifierare. \\
7 & \code{RX_DLC} & \code{0x1C} & R & Mottagen DLC. \\
8 & \code{RX_DATA_LO} & \code{0x20} & R & Mottagen data byte 0--3. \\
9 & \code{RX_DATA_HI} & \code{0x24} & R & Mottagen data byte 4--7. \\
10 & \code{RX_ACK} & \code{0x28} & W & Skriv \code{0x1} för att kvittera. \\
11 & \code{ERROR_FLAGS} & \code{0x2C} & R/W & Felregistret; skriv \code{0x0} för att nollställa. \\
12 & \code{TX_ABORT} & \code{0x30} & W & Skriv \code{0x1} för att tvinga tillbaka TX klar. \\
\end{booktable}

Index 13--15 är reserverade: en läsning ger \code{0x00000000}, en skrivning ignoreras.

\subsection{STATUS: de tre bitarna allt kretsar kring}

\begin{booktable}{@{}rlLL@{}}
\thead{Bit} & \thead{Betydelse} & \thead{Sätts av} & \thead{Nollställs av} \\ \midrule
0 & TX klar & reset; en avslutad sändning; en stigande flank på felnivån medan biten är låg; en
  skrivning till \code{TX_ABORT} & en accepterad skrivning till \code{TX_SEND} \\
1 & RX giltig & en mottagen frame, som samtidigt låser \code{RX_*} & en skrivning av \code{0x1}
  till \code{RX_ACK} \\
2 & Fel & en \textbf{stigande flank} på kontrollerns felnivå & en skrivning av \code{0x0} till
  \code{ERROR_FLAGS} \\
\end{booktable}

Bitarna 31--3 läses alltid som \code{0}.

\textbf{Läs tabellen en gång till.} Varje bit har en uttrycklig nollställning, och den
nollställningen är drivrutinens ansvar. Glöms \code{RX_ACK} står bit 1 kvar för alltid och du
läser samma frame i en oändlig loop. Glöms att \code{TX_SEND} nollställer bit 0 undrar du varför
den andra sändningen aldrig accepteras.

\subsection{Regler som är lätta att missa}

\begin{itemize}
  \item \code{TX_SEND} och \code{RX_ACK} kräver att \textbf{bit 0 är satt}; övriga bitar ignoreras.
  \item \code{ERROR_FLAGS} nollställs \textbf{bara av ett helt nollställt ord}. \code{0x1} gör
    ingenting, och \code{0x000000FF} gör ingenting trots att bit 0 är satt. Motsatt regel mot de
    två ovan.
  \item \code{TX_ID} maskerar till elva bitar, \code{TX_DLC} till fyra. En skrivning av
    \code{0xFFFFFFFF} läses tillbaka som \code{0x000007FF} respektive \code{0x0000000F}.
  \item Skrivningar till \code{STATUS} och \code{RX_*} ignoreras tyst.
  \item Läsningar av \code{TX_SEND} och \code{RX_ACK} ger \code{0x00000000}.
\end{itemize}

\subsection{Databyteordning}

\code{TX_DATA} och \code{RX_DATA} är 64 bitar breda, med \textbf{databyte 0 i de mest signifikanta
bitarna}:

\begin{console}
bit 63                                                              bit 0
+--------+--------+--------+--------+--------+--------+--------+--------+
| byte 0 | byte 1 | byte 2 | byte 3 | byte 4 | byte 5 | byte 6 | byte 7 |
+--------+--------+--------+--------+--------+--------+--------+--------+
 \___________ TX_DATA_LO ___________/ \___________ TX_DATA_HI __________/
\end{console}

\code{LO} och \code{HI} syftar på registerparets ordning, inte på bitsignifikans:
\textbf{\code{TX_DATA_LO} bär de första fyra databytesen.} Det är projektets vanligaste läsfel,
och det syns inte förrän en frame kommer fram bakvänd.

En frame med \code{dlc = 2} och data \code{AA BB} packas alltså som \code{TX_DATA_LO = 0xAABB0000}
och \code{TX_DATA_HI = 0x00000000}.

\subsection{De tre begränsningarna}
\label{sec:limits}

Hårdvaran är förenklad mot riktig CAN på många punkter. Tre av dem syns ända upp i din kod, och
måste hanteras snarare än döljas.

\begin{limitation}
\textbf{En förlorad arbitrering ser ut som en lyckad sändning.} \code{STATUS} bit 0 kommer tillbaka
i båda fallen. Biten betyder \emph{sändarporten är fri igen}, inte \emph{framen kom fram}. Vill du
veta utfallet: vänta på bit 0, och läs \textbf{sedan} \code{ERROR_FLAGS}.
\end{limitation}

\begin{limitation}
\textbf{\code{ERROR_FLAGS} är en bit för fyra orsaker} --- förlorad arbitrering, stoppfel,
misslyckad CRC och mottagen fjärrframe. Du kan se \emph{att} något gick fel, aldrig \emph{vad}.
Retry-logik som skiljer på dem går inte att skriva mot det här gränssnittet.
\end{limitation}

\begin{limitation}
\textbf{Ingen buffring på mottagarsidan.} Kontrollern håller exakt en frame. Kommer nästa innan
den förra kvitterats skrivs den över, och den nyaste vinner. Polla ofta --- och vet att det vid
någon trafikintensitet ändå inte räcker.
\end{limitation}

\begin{aside}
\note Hårdvaran validerar ingenting. En \code{TX_DLC} på \code{0xF} lagras och skickas vidare
precis som den är. Att avvisa ogiltiga frames är \textbf{drivrutinens} ansvar, och bara
drivrutinens: två kopior av samma regel är två saker som ska hållas i synk.
\end{aside}
